\section{Background}
One of the largest threats to global health in this century is antimicrobial resistance (AMR). It prolongs the duration of illnesses, increases the cost of treatment, and raises the chances of death. The World Health Organization (WHO) has issued a warning to countries to work together. Otherwise, the advantages of the antibiotic era may be lost and we will be heading back to a time where even a common infection may be fatal \cite{who_amr_2019}. Directly counting deaths from resistance is difficult. If a patient has a resistant infection and dies, the death is typically listed under the infection and not under the resistance. That is why researchers rely on statistical models built on hospital and mortality records. In 2019, an estimated 4.95 million deaths were associated with resistant bacterial infections, and 1.27 million of them were caused directly by the resistance \cite{murray2022gram}. Both figures are higher than the number of deaths from HIV/AIDS or from malaria. A second research group projects that 39 million people will die between 2025 and 2050 \cite{naghavi2024forecast}. These are model estimates rather than direct counts. The same trend is evident in hospital data. WHO has a programme called the Global Antimicrobial Resistance and Use Surveillance System, or GLASS. This programme enables laboratories in many countries to report antibiotic resistance of bacteria. The latest results show that about one bacterial infection in six does not respond to the antibiotic chosen for treatment. GLASS also monitors particular pairs of bacteria and antibiotics. Of these pairs, over 40\% were more difficult to treat from 2018 to 2023 \cite{who2025glass}. It's essential to make an early prediction of resistance. It helps physicians treat infections properly and helps prevent unnecessary broad-spectrum prescribing. At present, doctors rely on bacterial culture tests. These tests show how a strain responds to antibiotics, but they take 18 to 72 hours to complete. While waiting for results, a doctor may prescribe a broad-spectrum drug as a precaution, and this decision adds to the broader resistance problem. Some bacteria also cannot be cultured in a laboratory at all, and traditional approaches do not always identify the causes behind their resistance. As resistance continues to increase, tools that can predict resistance directly from sequence data are becoming more valuable.

\section{Rationale of the Study}
With the advance of whole genome sequencing, it is now possible to read the genetic sequence of a bacterial strain in detail. This has made it possible to predict resistance directly from DNA rather than waiting for a laboratory culture to grow. Many studies in machine learning have already reported high accuracy for antibiotic resistance prediction, which is discussed in detail in Chapter~\ref{ch:lit}. However, high accuracy on a benchmark does not necessarily mean that a model has learned the true biology of resistance.

An important factor underlying this limitation is bacterial population structure. Bacterial populations are strongly divided into lineages, and strains within the same lineage are likely to share their genetic background and resistance patterns. This issue is well documented in the literature on bacterial association testing. Because bacteria reproduce clonally, population structure and genome-wide linkage make it difficult to distinguish a mutation that truly causes resistance from one that is merely associated with it \cite{saber2020bacgwas}. Methods in that field now control for population structure directly within the model \cite{lees2020pyseer}, but this correction is almost never applied in the predictive machine learning literature. As a result, a model may appear to predict resistance very well simply by identifying the lineage a strain belongs to, without ever actually capturing the genes or mutations that underlie the resistance.

In practice, this distinction matters. A lineage-based model performs well on strains similar to its training examples, but it is far less reliable on a new strain. However, a new strain is also the one a clinician knows the least about, and it is exactly there that the model's answer carries the most weight. Predictions are least reliable where they are needed most. We are therefore motivated by two goals. First, we want a robust predictive model of antibiotic resistance in \textit{E. coli}, using genomic data as input. Second, and more importantly, we want to measure the source of that success: how much comes from true resistance mechanisms and how much comes from population structure. Only then can its predictions be trusted for the right reasons.

\section{Problem Statement}
Rapid and accurate prediction of antibiotic resistance remains a challenge, and this has not changed with modern sequencing and machine learning. However, there are still some issues that impact the doctor's choice. The first issue is time. It may take days to do a susceptibility test, and the doctor may have to use a broad-spectrum antibiotic until that test is done. When the chosen option is incorrect, the resistance may grow even worse. The second issue is coverage. The majority of tools assess resistance to individual drugs only. Doctors, however, often have to compare a number of drugs to select a treatment. The third issue is reach. Many tools are limited to searching for resistance genes that are known. Even if a strain is resistant, a tool may not provide a useful result if the strain does not contain any of these genes. These tools and their limitations are reviewed in Chapter~\ref{ch:lit}.

How these models are tested is just as important as how the models are built. Bacteria are grouped into families and strains in the same family can be highly similar. Close relatives can be found in the training set and test set during a random split. This means that the model can encounter a test strain similar to one it has already encountered in training. This can help to inflate the score. A more rigorous test employs strains that lack close relatives in the training data. But this kind of test is not provided in many studies.

Explanation is another issue. Many deep learning models give a prediction but do not explain why. Due to this fact, it is difficult to tell what the model has learned. It might have learned a gene or mutation that actually causes resistance. Or it may have learned just the family of the strain. Both would seem right from the outside of the model. But the first is true biology and the second is not. What is needed is a system that is faster, easier to understand, more thoroughly tested, and able to state the basis for its answers.

\medskip
\noindent \sloppy
This research addresses these concerns through the following research questions:
\begin{enumerate}
  \item[RQ1:] Whether a hybrid model that combines neural lenses with classical learners can predict antimicrobial resistance from genomic data.
  \item[RQ2:] How much each part of the model, the neural lenses, the gate, and the classical stacking layer, actually contributes.
  \item[RQ3:] Whether attention and gating weights can be used to interpret the model and identify which genes drive its predictions.
  \item[RQ4:] Whether the model's performance reflects real resistance mechanisms, population structure, or some mixture of the two.
  \item[RQ5:] Whether the model generalises to strains from lineages it has never seen.
\end{enumerate}

\section{Objective}
We aim to construct and validate a model to predict antibiotic resistance in \textit{E. coli} from genomic data, for seven antibiotics simultaneously. We also aim to quantify how much of its predictive capacity comes from true resistance mechanisms as opposed to population structure. Unlike a laboratory culture system, which takes a long time to produce a result, our system is designed to provide a fast prediction along with an honest account of what that prediction is based on. The work is structured in four steps to achieve this goal.

\medskip
\begin{itemize}
    \item Create a multi-antibiotic dataset controlled for leakage. We start from the publicly available Arcadia Science \textit{E. coli} genotype and phenotype resource, which combines approximately 7,000 strains from around the world with matching resistance labels \cite{morin2024ecoli}. A companion study on the same dataset \cite{york2025amr} examines the genetic basis of resistance and provides a biological framework for interpreting our results. For each genome, three views are constructed: SNP blocks clustered by gene, a gene presence/absence profile, and per-gene summary statistics. A separate collection of \textit{E. coli} from Giessen is kept aside. The complete data specification is given in Chapter~\ref{ch:method}.

    \item Develop TriLens-MoE-Stack. The three views are routed through three specialised neural lenses, which are fused through a per-sample, per-antibiotic mixture-of-experts gate, and then stacked with five classical base learners through a per-drug logistic regression meta-learner. Prediction for all seven antibiotics is done simultaneously as a multi-label task. In the gene presence and absence lens we discovered an issue with the standard gated attention and resolved it. This lens is based on a list of ones and zeros, which indicate the presence or absence of genes in a strain. With this type of input, the standard method gives the same score to every strain. So, it cannot properly separate one strain from another. We do not just say that our fix works. The benefit and the cost of it are measured by a controlled ablation.

    \item Evaluate generalisation on more than one kind of split. Results are shown for each of the following three splits: a random split, a phylogroup-stratified split, and a lineage holdout (based on sequence type of each strain). No sequence type is found in more than one partition. The second split is a stratification split and not a holdout. An overall report of all three makes the difference between the designs clear. We also rebuild and retrain the pipeline from scratch on the independent Giessen collection. This collection cannot support our three-view representation used for Arcadia, since it has only a single reference feature space, so this final test checks whether our methodology transfers to a new collection, not whether an already trained model transfers. The generalisation question could not be answered at all with a coarser split we tried first, as explained in Section~\ref{sec:comparisons}, and this failure is reported as a finding in itself.

    \item Measure lineage against mechanism. We combine the model's attention and gating weights with post-hoc attribution over the classical learners. The resulting marker genes are then functionally annotated and compared with a baseline given only general lineage information. Our intention was for the explanation to come from the model itself. In practice, the two sources do not fully agree, and we report that discrepancy rather than smoothing it over.
\end{itemize}

\section{Methodology in Brief}
This research is based on supervised learning of genomic data for \textit{E. coli}. After filtering, 6,962 isolates from the Arcadia collection are used together with an independent Giessen collection. Three feature views are generated for each genome, followed by a gene-filtering step that narrows the feature space down to the most informative genes. Resistance is represented as a multi-label task spanning seven antibiotics: AMP, CIP, AMC, SXT, GEN, CTX, and CAZ. Precision-recall AUC is the main performance measure, as it is appropriate for the class imbalance in resistance data. We also report clinical error rates and Matthews correlation. All scores include a bootstrap confidence interval. This allows the reader to spot if there is a meaningful difference between two scores or if that difference has arisen purely from noise. The random split is the reference that the rest are measured against. Harder tests follow: a phylogroup-stratified split, a sequence-type lineage holdout and a reconstruction of the whole pipeline on a second collection. One purposefully weak model is also trained. This model is only supplied with a broad family of the strain, no genome data. Its score indicates approximately the amount that is explained by family membership alone. The score of a real model should only be accepted if it is clearly better. The complete methodology is provided in Chapter~\ref{ch:method} and the results in Chapter~\ref{ch:results}.

\section{Scopes and Challenges}

\subsection{Scopes}
The work covers seven clinically important antibiotics and treats prediction as a multi-label problem, meaning one strain can be evaluated against all seven drugs at once, just as resistance occurs in practice. Interpretability is included as a core part of the work, not an add-on. There are various ways in which generalisation is tested. We take the splits that we have described above, and also reconstruct the entire method on a second collection. The framework is intended to be expandable to other species or other antibiotics in future studies, with possible applications in clinical diagnostics and public health surveillance. Clinical deployment, however, is specifically beyond the scope of this work. The system is a research prototype, not a clinical tool. Its very major error (VME) rate, the rate at which a truly resistant strain is reported as susceptible, is above the Clinical and Laboratory Standards Institute (CLSI) recommended error rate for all seven antibiotics evaluated. It is not designed to guide patient care.

\subsection{Challenges}
The first challenge was before the modelling stage. We required a large collection of genomes where the resistance labels were recorded in the same manner. Creating this data set was more challenging than anticipated. The raw data were in the form of variant call and PLINK files, which cannot be read directly by any modelling tool. Converting them on a single machine was slow, runs took hours, and the machine sometimes crashed. The tables we built from that data had their own problems, with many duplicate and overlapping columns, and none of them were grouped by gene. We therefore switched to a published variant release instead, as described in Section~\ref{sec:designadj}. The dataset remained large even after this change, occasionally causing processing failures. Memory was also a limiting factor, so intermediate results were saved throughout the analysis. Once training began, the classes were highly imbalanced, with far fewer resistant cases than susceptible ones. The models were able to predict susceptibility most of the time, so accuracy looked good, but this concealed the resistant cases that were missed. We therefore gave more weight to the rare class during training and stopped judging models by accuracy alone, using precision-recall AUC instead, since it focuses on detecting resistant cases. Consequently, there was a large variation in scores among the antibiotics tested. The patterns for amoxicillin-clavulanic acid were not generalisable to new data, and this was mitigated by early stopping and regularisation. We also kept a well-tuned classical model in the comparison, as a real competitor and not a formality.

\section{Thesis Organization}
\label{sec:organization}

The remainder of this report is structured such that each chapter responds to one question. This section provides a short map, so the reader can navigate directly to what they require.

Chapter~\ref{ch:lit} covers the background reading, and asks what is already known and what is not. It begins with the basic terms, such as what a bacterial genome is and how resistance is transferred from one strain to another. It then groups the earlier studies into six themes, and places this work next to the six closest of them. It closes by stating the gap that the rest of the report tries to close: how those models were tested, rather than how they were built.

Chapter~\ref{ch:req} sets out what the system had to do before any of it was built. It lists the functional and non-functional requirements, the limitations we had to operate under and the standards we followed. It also explores the social, environmental and ethical aspects of the work, together with the project plan, the risks and the cost. Fixing these targets in advance is what allows Chapter~\ref{ch:results} to report a requirement as unmet rather than quietly restate it.

Chapter~\ref{ch:method} describes the proposed system, and answers how it was built and on what data. The first half explains the design: three views of each genome, three neural lenses, a gate that weighs them, and a stacking layer on top of five classical learners. The second half describes the data, covering the origin of the genomes, how the labels were cleaned, how the features were reduced, and the manner in which the three evaluation splits were constructed. A reader who wants to repeat the experiments should begin here.

Chapter~\ref{ch:results} presents the results, and asks how much of the model's score survives a fair test. It begins with the random division, then measures the contribution of each component of the model. It then progresses to the more difficult tests: the lineage holdout, the prevalence correction the holdout requires, and the two-group pattern that appears once both are in place. It concludes by providing answers to the five research questions one after another. Chapter~\ref{ch:conclusion} draws the findings together, and asks what follows from them. It tells what the work says, where its limits are and what it should do next.
